\documentclass[11pt,twocolumn]{article}
\usepackage[margin=0.75in]{geometry}
\usepackage{booktabs}
\usepackage{amsmath}
\usepackage{amssymb}
\usepackage{graphicx}
\usepackage{hyperref}
\usepackage{xcolor}
\usepackage{listings}
\usepackage{microtype}
\usepackage{abstract}

% Colors and Styles
\definecolor{scryblue}{RGB}{30, 80, 160}
\hypersetup{colorlinks=true, linkcolor=scryblue, urlcolor=scryblue}

\lstset{
    basicstyle=\ttfamily\small,
    keywordstyle=\color{blue},
    stringstyle=\color{red},
    commentstyle=\color{gray},
    breaklines=true,
    frame=single,
    rulecolor=\color{lightgray}
}

\title{\textbf{Scry: A Unified Academic Suite for Pixel-Perfect Graphics and Machine Learning in Rust}}
\author{esoc \\ \texttt{scry-engine} workspace}
\date{February 2026}

\begin{document}

\twocolumn[
  \begin{center}
    {\LARGE \textbf{Scry: A Unified Academic Suite for Pixel-Perfect Graphics and Machine Learning in Rust}} \\
    \vspace{1em}
    {\large esoc} \\
    \textit{The scry-engine workspace} \\
    \vspace{1.5em}
  \end{center}
  \begin{onecolabstract}
    The Rust programming language has seen a surge in interest for scientific computing and systems programming, yet it lacks a unified, high-fidelity visualization and modeling ecosystem comparable to the Python scientific stack. We present \textbf{Scry}, a modular suite designed to bridge this gap. Scry provides a pixel-perfect graphics engine (\texttt{scry-engine}), a comprehensive charting library (\texttt{scry-chart}), a production-grade machine learning toolkit (\texttt{scry-learn}), and a feature engineering compiler (\texttt{scry-pipe}). By leveraging terminal graphics protocols (Kitty, Sixel) and GPU-accelerated rasterization, Scry enables publication-quality visualization directly within the developer's terminal, tightly integrated with model evaluation and data analysis.
  \end{onecolabstract}
  \vspace{2em}
]

\section{Introduction}
Scientific research and high-performance engineering require tools that are both fast and expressive. While Python dominates the data science landscape, the overhead of the Global Interpreter Lock (GIL) and the "two-language problem" often force developers to rewrite performance-critical logic in C++ or Rust. \texttt{Scry} is an attempt to provide a "single-language" solution for the entire lifecycle: from data preprocessing and modeling to high-fidelity visualization and production deployment.

\section{The Foundation: \texttt{scry-engine}}
At the core of the suite is \texttt{scry-engine}, a vector graphics engine specialized for modern terminal emulators. Unlike traditional TUI libraries that rely on Unicode block characters for "braille" graphics, \texttt{scry-engine} targets:
\begin{itemize}
    \item \textbf{Kitty Graphics Protocol:} Direct transmission of pixel data to the terminal.
    \item \textbf{Sixel:} Legacy-compatible raster graphics.
    \item \textbf{GPU Acceleration:} Integrated \texttt{wgpu} backend for rendering Signed Distance Fields (SDFs) and complex 3D scenes.
\end{itemize}

\section{Data Visualized: \texttt{scry-chart}}
\texttt{scry-chart} extends the engine into a high-level plotting library. It supports over 17 chart types, including:
\begin{itemize}
    \item \textbf{Standard:} Scatter, Line, Bar, Histogram.
    \item \textbf{Statistical:} BoxPlot, ViolinPlot, Heatmap.
    \item \textbf{Financial:} Candlestick, OHLC.
    \item \textbf{Specialized:} Radar, Waterfall, Funnel, and 3D Scatter.
\end{itemize}
Every chart is anti-aliased and supports multiple export targets (PNG, SVG, PDF), ensuring that terminal-based explorations can be directly used in academic papers.

\section{Intelligence Integrated: \texttt{scry-learn}}
\texttt{scry-learn} provides a robust implementation of classical machine learning algorithms. Unlike other Rust ML libraries, it emphasizes \textit{diagnostic integration}. 
\begin{lstlisting}[caption=Integrated Visualization in scry-learn]
let model = RandomForestClassifier::new();
model.fit(&train)?;
let report = classification_report(&test, &preds);
// Automatic high-fidelity chart generation
let chart = confusion_matrix_chart(&test, &preds);
\end{lstlisting}
The library covers anomaly detection, clustering (DBSCAN, K-Means), ensemble methods, and neural networks, all while maintaining strict numerical parity with established Python libraries.

\section{Deployment: \texttt{scry-pipe}}
The "deployment gap" is addressed by \texttt{scry-pipe}, a compiler that takes feature engineering pipelines defined in Python and generates zero-dependency Rust or WASM binaries. This ensures that scaling, imputation, and encoding logic remain identical across the training and serving environments.

\section{Performance and Benchmarks}
In preliminary audits, Scry has demonstrated significant performance advantages over Python-based alternatives in data-heavy TUI applications. By utilizing zero-copy transmission and shared memory (SHM) for the Kitty protocol, Scry achieves 60+ FPS for complex 3D visualizations within the terminal.

\section{Conclusion and Future Work}
Scry represents a paradigm shift in Rust-based scientific computing. By treating the terminal as a first-class graphical citizen and integrating modeling with visualization, it empowers researchers to work at the speed of thought without leaving their development environment. Future work includes expanding \texttt{scry-llm} into a full-featured tensor framework and broadening GPU-accelerated chart rendering.

\end{document}
